\begin{table}[htbp]
\centering
\caption{Robustness: 80th Percentile Dummy --- Excluding G-SIBs}
\label{tab:rob_p80_nocat1}
\small
\begin{tabular}{l c c c c c}
\toprule
 & (1) & (2) & (3) & (4) & (5) \\
 & Res. Mort. & Tot. Mort. & Oth. Mort. & Non-Mort. & NPL \\
\midrule
$\hat{\beta}_1$: High $\times$ Post 2020
 & -0.373 & -0.559 & -1.725 & 2.922 & 0.112 \\
 & (1.518) & (1.109) & (4.338) & (2.615) & (0.100) \\[4pt]
$\hat{\beta}_2$: High $\times$ Post 2022
 & -1.242 & -1.393 & 3.297 & 0.275 & 0.156 \\
 & (1.360) & (1.054) & (7.489) & (4.162) & (0.095) \\
\midrule
Observations   & 138 & 134 & 134 & 134 & 161 \\
$R^2$          & 0.314 & 0.401 & 0.245 & 0.413 & 0.807 \\
Bank FE        & \checkmark & \checkmark & \checkmark & \checkmark & \checkmark \\
Year FE        & \checkmark & \checkmark & \checkmark & \checkmark & \checkmark \\
\bottomrule
\end{tabular}
\begin{minipage}{\linewidth}
\vspace{4pt}
\footnotesize
\textit{Note:} Standard errors clustered at the bank level in
parentheses. CCyB exposure is standardised (mean 0, SD 1).
Column~(1): residential mortgage growth (direct CCyB target).
Column~(2): total mortgage growth.
Column~(3): other mortgage growth (within-mortgage spillover).
Column~(4): growth in non-mortgage loans (cross-segment spillover).
Column~(5): NPL ratio.
Binary treatment indicator based on 80th percentile split of CCyB exposure. Banks classified as high-exposure for the release (end-2019): Migros Bank AG, Bank Cler AG, Raiffeisen Group. For the reactivation (end-2021): Migros Bank AG, Luzerner Kantonalbank AG, Bank Cler AG. $\hat{\beta}_1$ (High $\times$ Post 2020) and $\hat{\beta}_2$ (High $\times$ Post 2022) capture the differential effect of being in the high-exposure group. Sample excludes Category~1 banks (UBS and Credit Suisse).
* \(p<0.1\), ** \(p<0.05\), *** \(p<0.01\)
\end{minipage}
\end{table}
